\begin{corollary}[折叠多重度剖面的中心奇矩全消失]\label{cor:derived-fold-centered-odd-moments-vanish}
沿用定理 \ref{thm:derived-fold-affine-reciprocity} 的记号。任取一个满足
\[
\widetilde\kappa_m\equiv \kappa_m\pmod{M_m}
\]
且与 \(M_m\) 奇偶相反的整数代表，并定义关于 \(\widetilde\kappa_m/2\) 的中心完全剩余系
\[
R_m:=
\left\{
r\in\ZZ:
\left|r-\frac{\widetilde\kappa_m}{2}\right|
<
\frac{M_m}{2}
\right\}.
\]
把 \(d_m\) 通过模 \(M_m\) 投影拉回到 \(R_m\) 后，对任意整数 \(j\ge 0\) 都有
\[
\sum_{r\in R_m}
d_m(r)
\left(r-\frac{\widetilde\kappa_m}{2}\right)^{2j+1}
=
0.
\]
等价地，中心指数型生成函数
\[
\Phi_m(t):=
\sum_{r\in R_m}
d_m(r)\,
e^{it(r-\widetilde\kappa_m/2)}
\]
\leanverified{paper\_derived\_fold\_centered\_odd\_moments\_vanish}
是 \(t\) 的偶函数；因而对每个 \(j\ge 0\) 都有
\[
\Phi_m^{(2j+1)}(0)=0.
\]
特别地，若在 \(t=0\) 邻域用 \(\log\Phi_m(t)\) 定义中心 cumulant 生成函数，则其全部奇阶 cumulant 都严格为零。对离散频率
\[
t=\frac{2\pi k}{M_m},
\qquad
k\in \ZZ/(M_m\ZZ),
\]
这与推论 \ref{cor:derived-fold-fourier-phase-rigidity} 中的实值性
\[
e^{\pi i k\kappa_m/M_m}\widehat d_m(k)\in \RR
\]
完全等价。
\end{corollary}

\begin{proof}
由于 \(\widetilde\kappa_m\) 与 \(M_m\) 奇偶相反，区间
\[
\left(
\frac{\widetilde\kappa_m-M_m}{2},
\frac{\widetilde\kappa_m+M_m}{2}
\right)
\]
的端点都是半整数，故其中恰含 \(M_m\) 个整数；这就给出 \(\ZZ/(M_m\ZZ)\) 的一个完全剩余系 \(R_m\)。同时，若 \(r\in R_m\)，则
\[
\left|
(\widetilde\kappa_m-r)-\frac{\widetilde\kappa_m}{2}
\right|
=
\left|
r-\frac{\widetilde\kappa_m}{2}
\right|
<
\frac{M_m}{2},
\]
故对合 \(r\mapsto \widetilde\kappa_m-r\) 保持 \(R_m\)。

再由定理 \ref{thm:derived-fold-affine-reciprocity}，
\[
d_m(r)=d_m(\widetilde\kappa_m-r)
\qquad
(r\in R_m),
\]
而
\[
\left(
(\widetilde\kappa_m-r)-\frac{\widetilde\kappa_m}{2}
\right)^{2j+1}
=
-\left(
r-\frac{\widetilde\kappa_m}{2}
\right)^{2j+1}.
\]
把 \(R_m\) 中元素按 \(r\leftrightarrow \widetilde\kappa_m-r\) 配对，相应项逐对相消，遂得所有中心奇矩都为零。

同理，
\[
\Phi_m(-t)
=
\sum_{r\in R_m}
d_m(r)\,
e^{-it(r-\widetilde\kappa_m/2)}
=
\sum_{r\in R_m}
d_m(\widetilde\kappa_m-r)\,
e^{it(r-\widetilde\kappa_m/2)}
=
\Phi_m(t),
\]
故 \(\Phi_m\) 为偶函数，从而 \(\Phi_m^{(2j+1)}(0)=0\)。又因
\[
\Phi_m(0)=\sum_{r\in R_m}d_m(r)>0,
\]
\(\log\Phi_m(t)\) 在 \(t=0\) 的某邻域内可取实解析分支，并同样是偶函数，于是其全部奇阶 Taylor 系数，即奇阶 cumulant，全都消失。最后，由
\[
\widetilde\kappa_m\equiv \kappa_m\pmod{M_m}
\]
可知离散频率处的相位校正与推论 \ref{cor:derived-fold-fourier-phase-rigidity} 完全一致。证毕。
\end{proof}
